\section{First-order face selection}\label{sec:firstorder}

In this section $T=\tau L$ with $\tau>0$ fixed, and we find the face that
carries the mode. For a face $F$ put
\[
 \Phi_\tau(F)=\tau\lambda_F+\lvert F\rvert-1 .
\]
We call $\Phi_\tau(F)$ the exponent of $F$. Theorem~\ref{thm:first} shows that
for every start the mode lies on an admissible face that minimizes
$\Phi_\tau$, and Theorem~\ref{thm:hull} describes these faces by a convex
hull. First we need the rate inside a face. For a face $F$ and $\alpha\in\Delta_F$ let
\[
 I_F(\alpha)=\sup_{u\in(0,\infty)^F}\sum_{i\in F}\alpha_i\Bigl(-\frac{(Q_Fu)_i}{u_i}\Bigr)
 =\sup_{u\in(0,\infty)^F}\sum_{i\in F}\alpha_i\Bigl(q_i-\sum_{j\in F\setminus\{i\}}q_{ij}\frac{u_j}{u_i}\Bigr),
\]
and let $q_i^{\rm out}(F)=\sum_{j\notin F}q_{ij}$. Taking $u=\ones$ gives
$I_F(\alpha)\ge\sum_i\alpha_iq_i^{\rm out}(F)$, and dropping the terms with
$q_{ij}$ gives $I_F(\alpha)\le\sum_i\alpha_iq_i$. As we said in the
introduction, $I_F$ is the Donsker--Varadhan rate of the occupation times of
the chain killed on leaving $F$. Dupuis and Liu~\cite{dupuisliu2015} give it
explicitly for reversible jump processes, and Guillin, Nectoux and
Wu~\cite{gnw2024} relate it to the rate of their conditional large deviation
principle when the process is reversible. We need its properties without
reversibility, and for a finite chain the proofs are short.

\begin{proposition}[the rate on a face]\label{prop:rate}
Let $F$ be irreducible with $k=\lvert F\rvert\ge2$.
\begin{enumerate}[(a)]
\item $I_F\ge\lambda_F$ on $\Delta_F$, with equality exactly at
$\alpha^*_F=l_F\circ r_F$.
\item $I_F$ is convex on $\Delta_F$ and strictly convex on $\Delta_F^\circ$. For
$\alpha\in\Delta_F^\circ$ the supremum defining $I_F(\alpha)$ is attained at a
vector $u=r_\alpha$, which is unique up to a positive factor, and $\alpha$ is
the stationary law of the generator on $F$ with off-diagonal entries
$q_{ij}r_{\alpha,j}/r_{\alpha,i}$.
\item If there is $\nu\in(0,\infty)^F$ with $\nu_iq_{ij}=\nu_jq_{ji}$ for all
$i,j\in F$, then for every $\alpha\in\Delta_F$
\[
 I_F(\alpha)=\sum_{\{i,j\}\subseteq F,\ i\ne j}\bigl(\sqrt{\alpha_iq_{ij}}-\sqrt{\alpha_jq_{ji}}\bigr)^2
 +\sum_{i\in F}\alpha_iq_i^{\rm out}(F).
\]
\item For every nonempty face $F'\subsetneq F$ and every
$\alpha\in\Delta_{F'}$, $I_F(\alpha)=I_{F'}(\alpha)\ge\lambda_{F'}>\lambda_F$.
\end{enumerate}
\end{proposition}

\begin{proof}
Write $u=e^\psi$, so that $I_F(\alpha)=\sup_{\psi\in\R^F}\Psi_\alpha(\psi)$
with
$\Psi_\alpha(\psi)=\sum_i\alpha_iq_i-\sum_{i\ne j}\alpha_iq_{ij}e^{\psi_j-\psi_i}$,
where $i,j$ run over $F$. Note that $\Psi_\alpha$ is concave in $\psi$, affine in
$\alpha$, and unchanged when a constant is added to $\psi$. Let $G^\psi$ be the
generator on $F$ with off-diagonal entries $q_{ij}e^{\psi_j-\psi_i}$. Then
$\nabla_\psi\Psi_\alpha(\psi)=-\alpha^\top G^\psi$, so $\psi$ maximizes
$\Psi_\alpha$ exactly when $\alpha$ is stationary for $G^\psi$. The switching
graph of $G^\psi$ is $G_F$, so $G^\psi$ is irreducible and has one stationary
law.

\emph{The lower bound in (a), for every face.} Let $F'$ be a nonempty face and
write $Q_{F'}=B-cI$ with $B\ge0$. For $s>\varrho(Q_{F'})$ we have
$s+c>\varrho(B)$, the spectral radius of $B$. So
$(sI-Q_{F'})^{-1}=\sum_{n\ge0}B^n(s+c)^{-n-1}\ge(s+c)^{-1}I$ entrywise, and
$u=(sI-Q_{F'})^{-1}\ones$ is positive with $Q_{F'}u=su-\ones<su$. This $u$
shows that $I_{F'}>-s$ on $\Delta_{F'}$. Letting $s\downarrow\varrho(Q_{F'})$
gives $I_{F'}\ge\lambda_{F'}$.

\emph{The value at $\alpha^*_F$.} Let $\psi=\log r_F$ and
$R=\operatorname{diag}(r_F)$. Since $Q_Fr_F=-\lambda_Fr_F$, we get
$G^\psi=R^{-1}(Q_F+\lambda_FI)R$, and then
$(\alpha^*_F)^\top G^\psi=l_F^\top(Q_F+\lambda_FI)R=0$. So $\log r_F$
maximizes $\Psi_{\alpha^*_F}$, and $I_F(\alpha^*_F)=\Psi_{\alpha^*_F}(\log
r_F)=\lambda_F$.

(b) Convexity holds because $I_F$ is a supremum of affine functions. Let
$\alpha\in\Delta_F^\circ$ and let $a$ be the least positive rate $q_{ij}$ with
$i,j\in F$. Normalize $\min_i\psi_i=0$ and put $M=\max_i\psi_i$. Some edge
$i\to j$ of a shortest path in $G_F$ from a minimum of $\psi$ to a maximum has
$\psi_j-\psi_i\ge M/(k-1)$, so
$\Psi_\alpha(\psi)\le\sum_i\alpha_iq_i-a(\min_i\alpha_i)e^{M/(k-1)}$, which
tends to $-\infty$ as $M\to\infty$. Hence the supremum is attained at some
$\psi$, and $\alpha$ is stationary for $G^\psi$.
If $\psi+\delta$ also attains it, then $\Psi_\alpha(\psi+t\delta)$ is constant
for $0\le t\le1$ by concavity. It is a constant minus a sum of convex functions
of $t$, so each term $\alpha_iq_{ij}e^{\psi_j-\psi_i+t(\delta_j-\delta_i)}$ is
affine in $t$. So $\delta_j=\delta_i$ for every edge $i\to j$ of $G_F$, and
$\delta$ is constant. For strict convexity let $\alpha^0\ne\alpha^1$ lie in
$\Delta_F^\circ$, let $\alpha^t=(1-t)\alpha^0+t\alpha^1$ with $0<t<1$, and let
$\psi$ attain $I_F(\alpha^t)$. Then
\[
 I_F(\alpha^t)=(1-t)\Psi_{\alpha^0}(\psi)+t\Psi_{\alpha^1}(\psi)\le(1-t)I_F(\alpha^0)+tI_F(\alpha^1),
\]
and equality would make both $\alpha^0$ and $\alpha^1$ stationary for $G^\psi$,
which is impossible.

(d) For $\alpha\in\Delta_{F'}$ only rows $i\in F'$ enter, and there the terms
with $j\in F\setminus F'$ are nonnegative and subtracted. Dropping them gives
$I_F(\alpha)\le I_{F'}(\alpha)$, and $u_j=\varepsilon\to0$ for $j\notin F'$
gives equality. The lower bound gives $I_{F'}(\alpha)\ge\lambda_{F'}$. Finally
$\lambda_{F'}>\lambda_F$, because the Perron root of an irreducible nonnegative
matrix, here $Q_F+cI$ for a large $c$, is larger than that of each proper
principal submatrix~\cite[Ch.~2]{bermanplemmons1994}.

(a) Let $\alpha\ne\alpha^*_F$. If $\alpha$ lies on the boundary of $\Delta_F$,
then (d) with $F'=\supp\alpha$ gives $I_F(\alpha)>\lambda_F$. If
$\alpha\in\Delta_F^\circ$ and $I_F(\alpha)=\lambda_F$, strict convexity gives
$I_F(\frac12(\alpha+\alpha^*_F))<\lambda_F$, which contradicts the lower bound.

(c) Since $q_i=\sum_{j\in F\setminus\{i\}}q_{ij}+q_i^{\rm out}(F)$, the
objective equals $\sum_i\alpha_iq_i^{\rm out}(F)$ plus a sum over pairs
$\{i,j\}$ of
$\alpha_iq_{ij}(1-u_j/u_i)+\alpha_jq_{ji}(1-u_i/u_j)$. By the AM--GM
inequality each such term is at most
$(\sqrt{\alpha_iq_{ij}}-\sqrt{\alpha_jq_{ji}})^2$. For the reverse inequality
put $u_i=\sqrt{\alpha_i/\nu_i}$ if $\alpha_i>0$ and $u_i=\varepsilon$ if
$\alpha_i=0$. If $\alpha_i,\alpha_j>0$, reversibility gives
$\alpha_iq_{ij}u_j/u_i=\alpha_jq_{ji}u_i/u_j=\sqrt{\alpha_i\alpha_jq_{ij}q_{ji}}$,
and the term equals its bound.
If $\alpha_i>0=\alpha_j$ the term tends to $\alpha_iq_{ij}$ as
$\varepsilon\to0$, and if $\alpha_i=\alpha_j=0$ both sides are $0$.
\end{proof}

When $Q_F$ is symmetric, $l_F=r_F$ and the minimizer $\alpha^*_F$ is the
entrywise square of the Perron vector normalized in $\ell^2$.

Away from $\alpha^*_F$ we need a bound that holds for all compositions on a
face at once. For an irreducible face $F$ put
$V_F=\{\theta\in\R^F:\varrho(A_F-\operatorname{diag}\theta)=0\}$. Note that
$q-\lambda_F\ones\in V_F$, since
$A_F-\operatorname{diag}(q-\lambda_F\ones)=Q_F+\lambda_FI$.

\begin{lemma}[a uniform upper bound]\label{lem:upper}
Let $F$ be irreducible with $k=\lvert F\rvert\ge2$, and let $\theta\in V_F$.
There is $C_\theta$, depending on $\theta$, $Q$ and $\mu$, such that for every
$n\in\Z_{\ge0}^F$ with $\supp n=F$ and every $h$ with $hq_{\max}\le\frac12$,
the probability that the first $\lvert n\rvert$ visits of the chain have
composition $n$ is at most
$C_\theta h^{k-1}\exp\bigl(-h\langle n,q-\theta\rangle+C_\theta h^2\lvert
n\rvert\bigr)$.
\end{lemma}

The proof is in Appendix~\ref{app:torus}. With $\theta=q-\lambda_F\ones$ it
gives $P_N(n)\le Ch^{k-1}e^{-T\lambda_F+CT^2/N}$ for every $n$ with support
$F$. The bound depends only on $n$ and $h$, and not on $N$, which we use for
the blocks of a path below.

\begin{theorem}[first-order face selection]\label{thm:first}
Fix $Q$ and $\mu$.
\begin{enumerate}[(i)]
\item Let $F$ satisfy (H), let $k=\lvert F\rvert$, and let
$\mathcal K\subset\Delta_F^\circ$ be compact. There is $N_0$ such that,
uniformly for $N\ge N_0$, $1\le T\le N^{1/3}$ and $\supp n=F$ with
$n/N\in\mathcal K$,
\[
 \log P_N(n)=-(k-1)L-T\,I_F(n/N)+O_{\mathcal K}(1+\log T).
\]
\item Let $F$ be admissible with components $C_1,\dots,C_r$, and let $k^*$ be
the size of a component $C_j$ with $\lambda_{C_j}=\lambda_F$. There are
$\kappa$, $T_0$ and $N_0$ such that for $N\ge N_0$ and $T_0\le T\le N^{1/3}$
\begin{align*}
 \log M_F&\ge-(k_F-1)L+(k_F-k^*)\log T-T\lambda_F-\kappa,\\
 \log M_F&\le-(k_F-1)L+(k_F-1)\log T-T\lambda_F+\kappa.
\end{align*}
If $F$ is not admissible, $M_F=0$.
\item Let $T=\tau L$ with $\tau>0$ fixed, and let $\mathcal F_\tau$ be the set
of admissible faces that minimize $\Phi_\tau$. Then
$\log\max_nP_N(n)=-L\min_F\Phi_\tau(F)+O(\log L)$, with the minimum over
admissible faces, and for large $N$ every global mode $\hat n$ has
$\supp\hat n\in\mathcal F_\tau$. This holds for every start $\mu$.
\item If in (iii) every face in $\mathcal F_\tau$ is irreducible, then every
global mode $\hat n$ satisfies $\hat n/N=\alpha^*_F+O(1/L)$ with
$F=\supp\hat n$. This is the case when $Q$ has symmetric support or $\mu$ has
full support (Theorem~\ref{thm:hull}(c)).
\end{enumerate}
\end{theorem}

Theorem~\ref{thm:local} and Lemma~\ref{lem:discrete} sharpen the error term
in (i) to $\frac{k-1}2\log T+\log C_F(n/N)+O(1/T+T^2/N)$, with $C_F$ from
Definition~\ref{def:tilt}.

\begin{remark}[the exponents and reducible faces]\label{rem:cost}
The factor $N^{-(k-1)}$ in (i) and (ii) comes from the $k-1$ switches that
enter new states (Remark~\ref{rem:lattice-term}), and $e^{-T\lambda_F}$ is
the chance of staying in $F$ for a time $T$. For a reducible face, the path in
our lower bound spends about $N/T$ visits, a time of order one, in every
component except one with the least exit rate. For such faces we only give bounds, and Proposition~\ref{prop:threecycle}
shows that the power of $T$ can differ from $(k_F-1)/2$.
\end{remark}

We prove (ii) and (iii) here. Parts (i) and (iv) follow from the second-order
results, Theorems~\ref{thm:local} and~\ref{thm:facemax} and
Lemma~\ref{lem:discrete}, and we derive them at the end of the proof.

\begin{proof}
(ii) If $F$ is not admissible, $M_F=0$ by definition. Let $F$ be admissible, let
$a$ be the least positive rate $q_{ij}$, and take $N_0$ so large that
$h\le N^{-2/3}$ is as small as we need. Note that $T^2/N\le1$. A path of
positive probability with support $F$ cannot leave a component and come back
to it, since the states visited in between would then belong to that
component. So it visits each component in one block of consecutive times and
passes from one block to the next by one switch. We call the order of the
blocks, the first state and these switches the pattern of the path. There are
finitely many patterns.

\emph{Upper bound.} Fix a pattern with blocks on
$C_{\pi(1)},\dots,C_{\pi(r)}$, first states $y_j$ and switches
$x_j\to y_{j+1}$, and let $\supp n=F$. Block $j$ has length
$N_j=\sum_{i\in C_{\pi(j)}}n_i$. The paths with this pattern and composition $n$
have total probability $\mu_{y_1}\prod_{j<r}hq_{x_jy_{j+1}}$ times the
product over $j$ of the total weight of the block paths from $y_j$ to $x_j$
inside $C=C_{\pi(j)}$ with the right composition. Without the condition on the last
state, this weight is the probability that the chain started at $y_j$ has that
composition in its first $N_j$ visits. By Lemma~\ref{lem:upper} with start $y_j$
and $\theta=q-\lambda_C\ones$ it is at most
$\kappa_1h^{\lvert C\rvert-1}e^{-h\lambda_CN_j+\kappa_1h^2N_j}$ when
$\lvert C\rvert\ge2$.
For $C=\{i\}$ it is $(1-hq_i)^{N_j-1}\le e^{1/2-h\lambda_CN_j}$. Since
$\sum_j(\lvert C_{\pi(j)}\rvert-1)+r-1=k_F-1$, $\sum_jN_j=N$ and
$\lambda_C\ge\lambda_F$ for every component, summing over the patterns gives
$P_N(n)\le\kappa_2(T/N)^{k_F-1}e^{-T\lambda_F+\kappa_1}$.

\emph{Lower bound.} By Lemma~\ref{lem:admissible}(a) we can order the
components as $C_1,\dots,C_r$ and choose $y_1\in C_1$ with $\mu_{y_1}>0$ and
$x_j\in C_j$, $y_{j+1}\in C_{j+1}$ with $q_{x_jy_{j+1}}>0$ for $j<r$. Take any
$x_r\in C_r$, let $\lambda_{C_{j^*}}=\lambda_F$ and $k^*=\lvert C_{j^*}\rvert$,
and put $m=\lceil N/T\rceil$, so that $1\le hm\le2$.

\emph{Claim A.} There is $c_A>0$ such that for every component $C$ and all
$y,x\in C$, the chain started at $y$ stays in $C$ for its first $m$ visits,
visits every state of $C$ and has $X_m=x$ with probability at least $c_A$.
Indeed, take a walk $y=v_0,\dots,v_J=x$ through every state of $C$ with
$q_{v_sv_{s+1}}>0$ and $J\le\lvert C\rvert^2$. The paths that make exactly these
switches, at any $J$ of the $m-1$ transition times, have total probability at
least $\binom{m-1}J(ha)^J(1-hq_{\max})^m\ge(1-Jh)^Ja^Je^{-4q_{\max}}/J!$,
since $1\le hm\le2$. Also $(1-Jh)^J\ge2^{-J}$ for small $h$.

\emph{Claim B.} There is $c_B>0$ such that for all $y,x\in C_{j^*}$ and
$\ell\ge2m$, the same event with $C_{j^*}$ and $\ell$ visits has probability at
least $c_Be^{-h\lambda_F\ell-h^2\lambda_F^2\ell}$. Indeed, let $w>0$ satisfy
$Q_{C_{j^*}}w=-\lambda_Fw$ and let $B=I+hQ_{C_{j^*}}\ge0$. Then $(B^s)_{zz'}$ is
the probability of going from $z$ to $z'$ in $s$ steps inside $C_{j^*}$, and
$\sum_{z'}(B^s)_{zz'}\ge(B^sw)_z/\max w\ge(1-h\lambda_F)^s\min w/\max w$. We
use Claim A for the first $m$ visits, from $y$ back to $y$, then $s=\ell-2m+1$
steps inside $C_{j^*}$, then Claim A for the last $m$ visits, ending at $x$.
This gives at least $c_A^2(\min w/\max w)(1-h\lambda_F)^\ell$. Also
$\lambda_F\le q_{\max}$, because the Perron root of $Q_{C_{j^*}}$ is at least
each diagonal entry. So $h\lambda_F\le\frac12$ and
$\log(1-h\lambda_F)\ge-h\lambda_F-h^2\lambda_F^2$.

Give block $j\ne j^*$ the length $m$ and block $j^*$ the length $N-(r-1)m$,
which is at least $2m$ for $T\ge T_0=r+2$ and $N\ge N_0$. The paths that go
from $y_j$ to $x_j$ in block $j$ as in the claims and pass from block $j$ to
block $j+1$ by the switch $x_j\to y_{j+1}$ all have support $F$. By the Markov
property their total probability is at least
$\mu_{y_1}(ha)^{r-1}c_A^{r-1}c_Be^{-T\lambda_F-q_{\max}^2}$. Their compositions
are fixed by the restrictions to the blocks, so there are at most
$(N+1)^{k^*-1}\prod_{j\ne j^*}(m+1)^{\lvert C_j\rvert-1}$ of them, and
$m+1\le3N/T$. Since $\sum_{j\ne j^*}(\lvert C_j\rvert-1)=k_F-k^*-(r-1)$,
dividing the probability by this number gives
$M_F\ge c\,N^{-(k_F-1)}T^{k_F-k^*}e^{-T\lambda_F}$.

(iii) Let $T=\tau L$. For large $N$ we have $T_0\le T\le N^{1/3}$, and by (ii)
$\log M_F=-L\,\Phi_\tau(F)+O(\log L)$ for every admissible face, while
$M_F=0$ for the others. Since $\max_nP_N(n)=\max_FM_F$ over finitely many
faces, the first claim follows. Let $m_\tau=\min_F\Phi_\tau(F)$ and let
$\varepsilon_0>0$ be the least value of $\Phi_\tau(F)-m_\tau$ over admissible
$F\notin\mathcal F_\tau$. Each such face has
$M_F\le N^{-m_\tau-\varepsilon_0+o(1)}$, which is smaller than
$\max_nP_N(n)=N^{-m_\tau+o(1)}$ for large $N$. A global mode has positive
probability, so its support is admissible, and hence it is in
$\mathcal F_\tau$.

\emph{Proof of (i) and (iv).} For (i) with $k=1$,
$P_N(Ne_i)=\mu_i(1-hq_i)^{N-1}$ with $\mu_i>0$ and $I_{\{i\}}=q_i$, so (i)
holds because $Nh^2\le1$. Let $k\ge2$, and let $f_T$ be the density of
Lemma~\ref{lem:continuum}. It is continuous and positive in $(T,\alpha)$, as
we note after Theorem~\ref{thm:local}. Also $C_F$ is positive and continuous on $\Delta_F^\circ$. So the
ratio of $f_T(\alpha)$ to $T^{(k-1)/2}C_F(\alpha)e^{-TI_F(\alpha)}$ is positive
and continuous, and by Theorem~\ref{thm:local} it tends to $1$ uniformly on
$\mathcal K$. Hence it is bounded above and below by positive constants for
$T\ge1$ and $\alpha\in\mathcal K$, and $\log f_T(\alpha)=-TI_F(\alpha)+O_{\mathcal
K}(1+\log T)$. By Lemma~\ref{lem:discrete},
$P_N(n)=N^{-(k-1)}f_T(n/N)(1+O(T^2/N))$, and the last factor lies in
$[\frac12,2]$ for $N\ge N_0$ since $T^2/N\le N^{-1/3}$. This proves (i).

For (iv), let $F\in\mathcal F_\tau$. It is admissible and irreducible, so
$\mu(F)>0$ by Lemma~\ref{lem:admissible}(b), and $F$ satisfies (H). By (iii),
for large $N$ a global mode $\hat n$ has its support $F$ in $\mathcal F_\tau$,
and it maximizes $P_N$ over the compositions with support $F$. Since
$T=\tau L\le N^{1/3}$ for large $N$, Theorem~\ref{thm:facemax} gives
$\hat n/N=\alpha^*_F+O(1/T)$, and there are finitely many faces.
\end{proof}

The law in (i) holds only on compact subsets of the open face, and the next
example shows that the boundary really is different.

\begin{example}[the boundary of a face]\label{ex:star}
Let $d=\ell+1$ with $\ell\ge3$, and let the only positive rates be $q_{jd}$ and
$q_{dj}$ for $j\le\ell$. So the switching graph is a star with center $d$ and
leaves $1,\dots,\ell$. Let $\mu$ have full support and $F=[d]$, so that $F$
satisfies (H). A leaf can only switch to the center, so the runs of a path
alternate between the center and the leaves. A path with support $F$ has at
least $\ell$ runs at leaves, and hence at least $\ell-1$ runs at the center.
So $P_N(n)=0$ whenever $\supp n=F$ and $n_d<\ell-1$, although $n/N$ lies in
$\Delta_F^\circ$ and $-(k-1)L-TI_F(n/N)$ is finite. Hence the law in
Theorem~\ref{thm:first}(i) cannot hold uniformly up to the boundary of
$\Delta_F$.
\end{example}

By Theorem~\ref{thm:first}(iii), first-order face selection is a finite
minimization. For an admissible face $F$ put
$P_F=(\lambda_F,\lvert F\rvert-1)\in\R^2$, and let $S$ be the set of these
points. Then $\Phi_\tau(F)=\tau x+y$ at $(x,y)=P_F$, so $\mathcal F_\tau$ is
the set of faces whose points minimize $\tau x+y$ on $S$. For each $\tau>0$
the points of $\operatorname{conv}S$ where $\tau x+y$ is least form a vertex
or an edge of $\operatorname{conv}S$. We call the union of these vertices and
edges over all $\tau>0$ the lower-left boundary of $\operatorname{conv}S$.

\begin{theorem}[geometry of the winners]\label{thm:hull}
Let $\mathcal F_\tau$ be the set of admissible faces that minimize
$\Phi_\tau$.
\begin{enumerate}[(a)]
\item The lower-left boundary of $\operatorname{conv}S$ is a path whose
vertices lie in $S$. Read from right to left, they are
$v_i=(\lambda^{(i)},y^{(i)})$, $0\le i\le m$, with
$\lambda^{(0)}>\dots>\lambda^{(m)}$ and $y^{(0)}<\dots<y^{(m)}$. Put
$\tau_i=(y^{(i)}-y^{(i-1)})/(\lambda^{(i-1)}-\lambda^{(i)})$ for $1\le i\le m$,
$\tau_0=0$ and $\tau_{m+1}=\infty$. Then $\tau_1<\dots<\tau_m$. For
$\tau_i<\tau<\tau_{i+1}$, $\mathcal F_\tau$ is the set of faces $F$ with
$P_F=v_i$. At $\tau=\tau_i$ it is the set of faces whose point lies on the
closed edge $[v_{i-1},v_i]$.
\item If $0<\tau<\tau'$, $F\in\mathcal F_\tau$ and $G\in\mathcal F_{\tau'}$,
then $\lambda_F\ge\lambda_G$ and $\lvert F\rvert\le\lvert G\rvert$. Both
inequalities are strict when $\mathcal F_\tau\cap\mathcal F_{\tau'}=\emptyset$.
\item If $Q$ has symmetric support, or $\mu$ has full support, then every face
in $\mathcal F_\tau$ is irreducible, for every $\tau>0$.
\end{enumerate}
\end{theorem}

\begin{proof}
(a) Let $g(\tau)=\min_{(x,y)\in S}(\tau x+y)$ for $\tau>0$. It is a minimum of
finitely many affine functions, so it is concave and piecewise affine. Let
$0<\tau_1<\dots<\tau_m$ be the points where its slope changes, with
$\tau_0=0$ and $\tau_{m+1}=\infty$. On $(\tau_i,\tau_{i+1})$ we have
$g(\tau)=\tau\lambda^{(i)}+y^{(i)}$ for some
$v_i=(\lambda^{(i)},y^{(i)})\in S$. If $p=(x,y)\in S$ has $\tau x+y=g(\tau)$
for some $\tau$ in this interval, then $\tau'\mapsto\tau'x+y-g(\tau')$ is
affine and nonnegative on the interval and vanishes inside it. So it vanishes
on the interval and $p=v_i$. The slope of $g$ decreases, so
$\lambda^{(0)}>\dots>\lambda^{(m)}$, and continuity at $\tau_i$ gives
$\tau_i\lambda^{(i-1)}+y^{(i-1)}=\tau_i\lambda^{(i)}+y^{(i)}$. This is the
formula for $\tau_i$, and it gives $y^{(i)}>y^{(i-1)}$. Now let $p=(x,y)\in S$
with $\tau_ix+y=g(\tau_i)$. It lies on the line $\tau_ix+y=g(\tau_i)$, which
contains $v_{i-1}$ and $v_i$. Comparing $\tau x+y$ with $g(\tau)$ just below
and just above $\tau_i$ gives $\lambda^{(i)}\le x\le\lambda^{(i-1)}$, so $p$
lies on the closed edge $[v_{i-1},v_i]$. Conversely every point of $S$ on
this edge attains $g(\tau_i)$. For each $\tau$ the minimizers of $\tau x+y$ on
$\operatorname{conv}S$ form the convex hull of its minimizers on $S$, which is
$v_i$ or $[v_{i-1},v_i]$. So the lower-left boundary is the path
$v_0v_1\cdots v_m$. Its edges are exposed by different values of $\tau$, so
they have different directions and the $v_i$ are its vertices.

(b) By optimality $\Phi_\tau(F)\le\Phi_\tau(G)$ and
$\Phi_{\tau'}(G)\le\Phi_{\tau'}(F)$. Adding them gives
$(\tau'-\tau)(\lambda_F-\lambda_G)\ge0$, so $\lambda_F\ge\lambda_G$, and then
the first inequality gives
$\lvert F\rvert-\lvert G\rvert\le\tau(\lambda_G-\lambda_F)\le0$. If the
sets are disjoint and $\lambda_F=\lambda_G$, the inequalities give
$\lvert F\rvert=\lvert G\rvert$, so $\Phi_{\tau'}(F)=\Phi_{\tau'}(G)$ and
$F\in\mathcal F_{\tau'}$, which is a contradiction. So $\lambda_F>\lambda_G$,
and then $\lvert F\rvert<\lvert G\rvert$.

(c) With symmetric support every admissible face is irreducible by
Lemma~\ref{lem:admissible}(c). Let $\mu$ have full support, and let $F$ be
admissible and reducible. Take a component $C$ of $F$ with
$\lambda_C=\lambda_F$. It is irreducible with $\mu(C)>0$, so it is admissible
by Lemma~\ref{lem:admissible}(b). Since $\lvert C\rvert\le\lvert F\rvert-1$,
we get $\Phi_\tau(C)\le\Phi_\tau(F)-1$, so $F$ is not a minimizer.
\end{proof}

A face with small $\lambda_F$ for its size is a set of states that the chain
rarely leaves, a well-separated community. As $\tau$ grows the exit rate
counts for more against the size, and the mode moves to larger communities
that hold the chain better. It need not grow around its first vertex, and it
can move to a different community (Remark~\ref{rem:nonnested}).
Figure~\ref{fig:cycles}(a) shows $S$ and its lower-left boundary for the
6-cycle.

\begin{figure}[t]
\centering
\includegraphics{fig_cycles}
\caption{(a) The points $(\lambda_F,\lvert F\rvert-1)$ of the connected faces
of the 6-cycle $C_6$. These are the arcs of $1$ to $5$ states and $C_6$
itself, and all arcs of one size give the same point. The line is the
lower-left boundary of the convex hull. The winners are the arcs of $1$ to $4$
states and then $C_6$ (filled), while the arc of $5$ states never wins (open).
The switch values $1$, $1+\sqrt2$, $4.906$ and $3+\sqrt5$ are minus the slopes
of the hull edges. (b) The support of the exact mode on $C_5$ with the uniform
start, against $T$, at $N=40$ and $N=72$. The digits give the number of states
in the support. The mode changes at $T=2.538$, $5.123$ and $6.884$ for $N=40$
and at $T=3.078$, $6.389$ and $8.602$ for $N=72$. The triangles mark the
first-order switch points $\tau\log N$ with $\tau=1$, $1+\sqrt2$ and
$2+\sqrt2$ (Theorem~\ref{thm:cycle}).}
\label{fig:cycles}
\end{figure}

For symmetric generators the exit rates have a variational form. For
$1\le s\le d$ let
$\underline\lambda_s=\min\{\lambda_F:F\text{ admissible},\ \lvert F\rvert=s\}$
when such a face exists. A face of size $s$ with
$\lambda_F>\underline\lambda_s$ has a larger value of $\Phi_\tau$ than one
with $\lambda_F=\underline\lambda_s$, for every $\tau>0$, so only the points
$(\underline\lambda_s,s-1)$ matter. In part (c) of the next lemma we plot them
with the axes swapped, as $(s-1,\underline\lambda_s)$, so that the winners are
read off a lower convex hull.
We call a face connected if its switching graph is connected.

\begin{lemma}[symmetric generators]\label{lem:symmetric}
Let $Q$ be symmetric and irreducible, and put
$\mathcal E(f)=\sum_{i<j}q_{ij}(f_i-f_j)^2$ for $f\in\R^d$.
\begin{enumerate}[(a)]
\item For every face $F$,
$\lambda_F=\min\{\mathcal E(f):\supp f\subseteq F,\ \lVert f\rVert_2=1\}$, and
the minimum is attained at some $f\ge0$.
\item If $F\subsetneq G$ and $G$ is connected, then $\lambda_G<\lambda_F$.
\item If $\mu$ has full support, then
$\underline\lambda_1>\underline\lambda_2>\dots>\underline\lambda_d=0$. Let
$1=s_0<s_1<\dots<s_m=d$ be the sizes $s$ at which $(s-1,\underline\lambda_s)$
is a vertex of the lower convex hull of these $d$ points, and put
$\tau_i=(s_i-s_{i-1})/(\underline\lambda_{s_{i-1}}-\underline\lambda_{s_i})$
for $1\le i\le m$, $\tau_0=0$ and $\tau_{m+1}=\infty$. Then $\tau_1<\dots<\tau_m$, and for
$\tau_i<\tau<\tau_{i+1}$ the set $\mathcal F_\tau$ consists of the faces $F$
with $\lvert F\rvert=s_i$ and $\lambda_F=\underline\lambda_{s_i}$.
\end{enumerate}
\end{lemma}

\begin{proof}
(a) Since $Q$ is symmetric with zero row sums, $\langle
f,-Qf\rangle=\mathcal E(f)$. If $\supp f\subseteq F$ this equals $\langle
f_F,-Q_Ff_F\rangle$, and the least eigenvalue of the symmetric matrix $-Q_F$ is
$\lambda_F$. So (a) is the Rayleigh--Ritz principle, together with $\mathcal
E(\lvert f\rvert)\le\mathcal E(f)$.

(b) Monotonicity follows from (a). Let $\varphi\ge0$ attain the minimum for
$F$, and suppose that $\lambda_G=\lambda_F$. Then $\varphi$ also minimizes the
Rayleigh quotient of $-Q_G$, so $-Q_G\varphi=\lambda_G\varphi$ on $G$. Since
$\supp\varphi\subseteq F\ne G$ and $G$ is connected, some
$j\in G\setminus\supp\varphi$ has $q_{ji}>0$ for an $i\in\supp\varphi$. Then
$(-Q_G\varphi)_j=-\sum_{i'}q_{ji'}\varphi_{i'}<0=\lambda_G\varphi_j$, which is
a contradiction.

(c) Here the admissible faces are the connected faces, by
Lemma~\ref{lem:admissible}(b),(c). Also $\lambda_{[d]}=0$ because $Q\ones=0$
and $\ones>0$, so $\underline\lambda_d=0$. Let $s<d$ and let $F$ be connected
with $\lvert F\rvert=s$ and $\lambda_F=\underline\lambda_s$. The switching
graph of $Q$ is connected, so some $j\notin F$ has a neighbor in $F$. Then
$F\cup\{j\}$ is connected, and (b) gives
$\underline\lambda_{s+1}\le\lambda_{F\cup\{j\}}<\underline\lambda_s$. The
lower convex hull of the points $(s-1,\underline\lambda_s)$ is the graph of a
convex piecewise linear function with corners at the points
$(s_i-1,\underline\lambda_{s_i})$. Its slopes
$\sigma_i=(\underline\lambda_{s_i}-\underline\lambda_{s_{i-1}})/(s_i-s_{i-1})
=-1/\tau_i$ increase with $i$, and they are negative because
$\underline\lambda_s$ decreases. Minimizing $\tau\underline\lambda_s+s-1$
over $s$ means finding where a line of slope $-1/\tau<0$ supports this graph.
With $\sigma_0=-\infty$ and $\sigma_{m+1}=0$, the line touches only the corner
at $s_i$ exactly when $\sigma_i<-1/\tau<\sigma_{i+1}$, that is when
$\tau_i<\tau<\tau_{i+1}$.
\end{proof}
